\documentclass[10pt,a4paper]{article}

% ---------- Seiten- und Schrifteinstellungen ----------
\usepackage[a4paper,top=1.20cm,bottom=1.20cm,left=1.55cm,right=1.55cm]{geometry}
\usepackage[T1]{fontenc}
\usepackage[utf8]{inputenc}
\usepackage[ngerman]{babel}
\usepackage[scaled=0.96]{helvet}
\renewcommand{\familydefault}{\sfdefault}
\usepackage{microtype}
\usepackage{graphicx}
\usepackage{tabularx}
\usepackage{enumitem}
\usepackage{xcolor}
\usepackage{titlesec}
\usepackage[hidelinks]{hyperref}
\usepackage{ragged2e}
\usepackage{fancyhdr}
\usepackage{lastpage}

% ---------- PDF-Metadaten ----------
\hypersetup{
  pdftitle={REPLACENAME - Lebenslauf Maschinenbau},
  pdfauthor={REPLACENAME},
  pdfsubject={Maschinenbau, Produktentwicklung, CAE und Strukturdynamik}
}

% ---------- Stil ----------
\definecolor{heading}{HTML}{202020}
\definecolor{muted}{HTML}{555555}
\definecolor{rulecolor}{HTML}{777777}

\pagestyle{fancy}
\fancyhf{}
\renewcommand{\headrulewidth}{0pt}
\renewcommand{\footrulewidth}{0pt}
\fancyfoot[L]{\scriptsize\color{muted}REPLACENAME}
\fancyfoot[R]{\scriptsize\color{muted}\thepage/\pageref{LastPage}}

\setlength{\parindent}{0pt}
\setlength{\parskip}{0pt}
\setlist[itemize]{leftmargin=1.18em,itemsep=2.0pt,topsep=2.5pt,parsep=0pt,partopsep=0pt}
\urlstyle{same}
\RaggedRight

\titleformat{\section}
  {\large\bfseries\color{heading}}
  {}{0pt}{}
  [\vspace{-0.22em}\color{rulecolor}\titlerule]
\titlespacing*{\section}{0pt}{0.55em}{0.30em}

% Einträge über die volle Breite: Inhalt links, Datum bzw. Angabe rechts.
\newcommand{\cvjob}[4]{%
  \begin{tabularx}{\linewidth}{@{}Xr@{}}
    \textbf{#1} & \textbf{#2}\\[-1pt]
    \textit{#3} & \textit{#4}
  \end{tabularx}
  \vspace{0.05em}
}

\newcommand{\cvproject}[2]{%
  \begin{tabularx}{\linewidth}{@{}Xr@{}}
    \textbf{#1} & \textbf{#2}
  \end{tabularx}
  \vspace{0.05em}
}

\newcommand{\cvedu}[4]{%
  \noindent
  \textbf{#1}\hfill\mbox{\textbf{#2}}\par
  \vspace{2pt}
  \noindent
  \textit{#3}\hfill\mbox{\textit{#4}}\par
  \vspace{0.03em}
}
\hyphenation{SOLIDWORKS OpenFAST}
\emergencystretch=1em
\renewcommand{\arraystretch}{1.08}
\begin{document}
\fontsize{11.25}{12.75}\selectfont

% ---------- Kopfbereich ----------
\begin{minipage}[t]{0.79\linewidth}
  \vspace{0pt}
  {\fontsize{21.5}{23.5}\selectfont\bfseries REPLACENAME}\\[2pt]
  {\fontsize{11.55}{13.2}\selectfont\bfseries Maschinenbauingenieur | Produktentwicklung, CAE \& Strukturdynamik}\\[4pt]
  REPLACELOCATION \enspace|\enspace
  \href{tel:REPLACEPHONE}{REPLACEPHONE} \enspace|\enspace
  \href{mailto:REPLACEEMAIL}{REPLACEEMAIL}\\[1.5pt]
  \href{REPLACELINKEDIN}{REPLACELINKEDIN} \enspace|\enspace
  \href{REPLACEPORTFOLIO}{REPLACEPORTFOLIO}\\[1.5pt]
\end{minipage}
\hfill
\begin{minipage}[t]{0.17\linewidth}
  \vspace{0pt}
  \raggedleft
  \IfFileExists{candidate_photo.jpeg}{\includegraphics[width=2.35cm,height=3.00cm,keepaspectratio]{candidate_photo.jpeg}}{\rule{0pt}{3.00cm}\hspace{2.35cm}}
\end{minipage}

\section{Kurzprofil}
Maschinenbauingenieur mit nahezu sechs Jahren Industrieerfahrung in der Entwicklung von Maschinen, Mechanismen für erneuerbare Energien und mechanischen Systemen -- von Konzeption und Analyse bis zu Fertigung und Installation. Erfahrung in 3D-CAD, statischen und dynamischen CAE-Analysen, der Entwicklung von Schwingungssystemen, Komponentenauswahl, Fertigungsdokumentation und bereichsübergreifender Produktionsunterstützung. Derzeitiger Schwerpunkt im M.Sc.-Studium Wind Energy Engineering: Anwendung fundierter Kenntnisse in mechanischer Simulation und Engineering-Automatisierung auf Lastanalysen von Windenergieanlagen, aeroelastische Simulationen und Windpark-Standortbewertungen.

\section{Technische Kernkompetenzen}
\begin{tabularx}{\linewidth}{@{}p{5.6cm}X@{}}
\textbf{Produktentwicklung} & Anforderungsdefinition, Konzeptentwicklung, Detailkonstruktion, fertigungsgerechte Konstruktion, Unterstützung bei der Installation\\
\textbf{CAD \& Dokumentation} & CATIA V5, SOLIDWORKS, Siemens NX, Fertigungszeichnungen, Stücklisten, SAP-basierte Produktdaten\\
\textbf{CAE \& Dynamik} & ANSYS Workbench, Strukturanalyse, Modalanalyse, harmonische Analyse, Schwingungs- und Resonanzbewertung\\
\textbf{Maschinen \& Mechanismen} & Schwingmaschinen, Auswahl von Motoren und Erregersystemen, Solar-Nachführung und automatisierte Reinigungsmechanismen\\
\textbf{Programmierung \& Simulation} & Python, MATLAB/Simulink, OpenFAST, Power BI, Automatisierung technischer Arbeitsabläufe\\
\textbf{Projektumsetzung} & Fertigungs- und Montageunterstützung, technische Koordination, Konstruktionsänderungen, Kommunikation mit Lieferanten und Fertigung
\end{tabularx}

\section{Berufserfahrung}
\cvjob{Entwicklungsingenieur - Solar-Trackingsysteme \& automatisierte Mechanismen}{01/2022--02/2023}{Paydar | Solarenergiesysteme}{Teheran}
\begin{itemize}
  \item Konstruktion von Photovoltaik-Nachführsystemen, Befestigungskomponenten und mechanischen Baugruppen in SOLIDWORKS für eine zuverlässige Fertigung und Installation.
  \item Optimierung von Komponenten mithilfe struktureller Analysen in ANSYS; Beitrag zu einer Reduzierung von Produktionsausschuss und Rohmaterialabfall um etwa 10\%.
  \item Entwicklung eines automatisierten Reinigungssystems für Solarmodule mit Elektromotoren, Führungsschienen und Rädern für eine kontrollierte Bewegung über die Modulflächen.
  \item Umsetzung des Reinigungssystems von der Konzept- und Detailkonstruktion bis zu Fertigung und Montage.
  \item Erstellung von Fertigungszeichnungen, Stücklisten und Komponentenspezifikationen sowie Pflege technischer Daten in SAP-basierten Systemen.
\end{itemize}

\cvjob{Maschinenbauingenieur - Maschinenentwicklung \& Strukturdynamik}{03/2017--12/2021}{TehranRigzar | Bergbau- \& Aufbereitungstechnik}{Teheran}
\begin{itemize}
  \item Konstruktion mechanischer Anlagen und Baugruppen für die Gesteinsaufbereitung, darunter Backenbrecher, Förderbänder, Hochfrequenz-Schwingsiebe und Waschanlagen für Zuschlagstoffe.

  \item Erstellung von Maschinenbaugruppen, Strukturbauteilen, mechanischen Schnittstellen und Fertigungszeichnungen in CATIA V5.

  \item Mitentwicklung eines Hochfrequenz-Schwingsiebs auf Grundlage von Laboruntersuchungen zur Siebung und definierten Betriebsanforderungen.

  \item Mitwirkung beim Ersatz des bisherigen kurbelwellengetriebenen Siebs durch eine Konstruktion mit Vibrationsmotoren zur Verbesserung der Zuverlässigkeit und Vereinfachung der Wartung.

  \item Durchführung von Modal- und harmonischen Analysen in ANSYS zur Bewertung des dynamischen Verhaltens des Siebsystems.

  \item Zusammenarbeit mit dem Entwicklungsteam bei Betriebsberechnungen und der Auswahl der Vibrationsmotoren für die endgültige Maschinenkonfiguration.

\end{itemize}

\newpage

\section{Aktuelle technische Spezialisierung}
\cvproject{M.Sc. Wind Energy Engineering - Hochschule Flensburg}{03/2023--heute}
\begin{itemize}
  \item \textbf{Masterarbeit:} Entwicklung eines Rotor-als-Sensor-Ansatzes zur Schätzung von Windgeschwindigkeit, Windrichtung, vertikaler Windscherung und Winddrehung aus messbaren Anlagensignalen.

  \item Automatisierung der Simulations- und Signalverarbeitungsabläufe in Python mit geplantem Vergleich linearer Schätzer und neuronaler Netze.

  \item Analyse aeroelastischer Lasten und der Strukturreaktion einer 20-MW-Offshore-Windenergieanlage mit Monopile-Gründung mithilfe von OpenFAST, Python und MATLAB.

  \item Durchführung einer Windparkplanung und Umweltbewertung in windPRO und QGIS mit Anlagenlayout, Energieertrag, Schallimmissionen, Schattenwurf und Abschaltungen zum Fledermausschutz.
\end{itemize}

\section{Ausgewählte Ingenieurprojekte}
\cvproject{Autonomer Unterwasser-Schweißroboter - Bewegungssteuerung}{M.Sc.-Abschlussprojekt}
\begin{itemize}
  \item Entwicklung und Implementierung eines Bewegungssteuerungssystems in MATLAB/Simulink zur präzisen Bewegung und Positionierung eines autonomen Unterwasser-Schweißroboters.
\end{itemize}

\cvproject{Handprothese \& additive Fertigung}{Ehrenamtliches Projekt}
\begin{itemize}
  \item Konstruktion und Prototypenfertigung mechanischer Komponenten für eine Handprothese mittels 3D-Druck; Unterstützung der iterativen Entwicklung vom digitalen Modell bis zum gefertigten Bauteil.
  \item Arbeit mit 3D-Druckerhardware und Fertigungsparametern zur Herstellung und Optimierung der Prototypenbauteile.
\end{itemize}

\section{Studium}
\cvedu
  {M.Sc. Wind Energy Engineering -- Notenschnitt: 1,9}
  {03/2023--heute}
  {Hochschule Flensburg, Flensburg, Deutschland}
  {Abschluss: Februar 2027}

\cvedu{M.Sc. Mechanical Engineering}{09/2014--01/2017}{Hakim Sabzevari University, Sabzevar}{Regelungstechnik und Robotik}

\cvedu{B.Sc. Mechanical Engineering}{09/2008--06/2013}{Azad University, Teheran}{Maschinenkonstruktion und Schwingungsanalyse}

\section{Ausgewählte Weiterbildungen \& Zertifikate}
\begin{tabularx}{\linewidth}{@{}p{4.5cm}X@{}}
\textbf{Python-Programmierung} & HarvardX\\
\textbf{ANSYS} & IEEE-Schulung\\
\textbf{SOLIDWORKS} & Technical and Vocational Training Organization\\
\textbf{CATIA V5} & Amirkabir University of Technology
\end{tabularx}

\section{Sprachen}
\begin{tabularx}{\linewidth}{@{}p{4.5cm}X@{}}
\textbf{Englisch} & C1; IELTS-Gesamtergebnis 7,0\\
\textbf{Deutsch} & B1; Vorbereitung auf eine Zertifikatsprüfung
\end{tabularx}

\end{document}
